\documentclass[10pt,a4paper]{amsart}
\usepackage[margin=19mm]{geometry}
\usepackage{amsmath,amssymb,mathtools}
\usepackage[scr=rsfs,cal=euler]{mathalfa}
\usepackage{xfrac}
\usepackage[unicode,hidelinks,bookmarks=false]{hyperref}
\newcommand{\ie}{i.e.\ }
\newcommand{\cf}{cf.\ }
\newcommand{\resp}{respectively\ }
\newcommand{\Subsection}{\subsection}
\providecommand{\nobreakdash}{\nobreak\hbox{-}}
\setlength{\emergencystretch}{2em}
\multlinegap0pt
\usepackage{mathtools}
\usepackage{amssymb}
\newtheorem{theorem}{Theorem}
\title{Orthogonal Latin Squares of Order Ten with Two Relations: A SAT Investigation}
\author{Curtis Bright, Amadou Keita, Brett Stevens}
\date{Source: arXiv:2509.09633v3 (CC BY 4.0)}
\begin{document}
\maketitle

\noindent\emph{Source and licence.} Orthogonal Latin Squares of Order Ten with Two Relations: A SAT Investigation. Version arXiv:2509.09633v3 (CC BY 4.0), distributed by arXiv under the Creative Commons Attribution 4.0 International licence, http://creativecommons.org/licenses/by/4.0/, as recorded in the arXiv licence metadata for this exact version and shown on the abstract page https://arxiv.org/abs/2509.09633v3. Full licence notice: \texttt{licenses/arxiv-2509.09633-CC-BY-4.0.txt}.\par\smallskip

\noindent\textit{Editorial scope.} Counted unit: the article's theorem thm:pointtype (source0.tex lines 757--759). The article's complete original proof of it is delivered with it (source0.tex lines 761--793, one \texttt{proof} environment of 33 lines). No mathematics is written, repaired or re-derived: the statement and the proof are the article's own lines, in the article's own notation, and no retained line is substituted or re-flowed.  Retained uncounted as the source-local context the proof needs: lines 147--147.  Nothing was omitted from any retained span, so the package carries no omission record.  No retained line contains a citation, so the wrapper carries no bibliography and no bibliography entry.  No retained line contains a figure environment, an image-inclusion command or drawing code, so no image is delivered and the package declares no asset dependency.  Editorial formatting only, in addition: an AMS \texttt{amsart} preamble that restates the article's own declarations and macros used by the retained lines, namely the environments \texttt{theorem} on separate counters, together with the standard packages those lines need.  The retained environments print in this document as Theorem 1 in order of appearance, whereas the article prints the same items with the numbering of its own full document.  The complete native source of the article as distributed in the arXiv e-print for this exact version is bundled unchanged as \texttt{sources/184922/source0.tex}.
\section{Background}\label{sec:background}

\begin{theorem}\label{thm:pointtype}
There exist no 4-nets(10) with two relations in cases~2--4.
\end{theorem}

\begin{proof}
Say $R_1$ and $R_2$ are two relations in a 4-net(10).
In what follows we use the relational code `0' to denote $R_1\cap R_2$, `1' to denote $R_1\setminus R_2$,
`2' to denote $R_2\setminus R_1$, and `3' to denote $\overline{R_1}\cap\overline{R_2}$.
The \emph{type} of a point is a 4-character $\{0,1,2,3\}$-string denoting the point's relational codes
from each parallel class.  For example, a point of type 1122 is in $R_1$ but not $R_2$ in the first
two classes, and is in $R_2$ but not $R_1$ in the last two classes.
Let $t_{ijkl}$ denote the number of points in the net of type $ijkl$.
Note that for $R_2$ to be a relation it must be the case that $t_{ijkl}=0$ when $i+j+k+l\not\equiv0\pmod{2}$,
and similarly for $R_1$ to be a relation it must be the case that $t_{ijkl}=0$ when
$\lfloor i/2\rfloor+\lfloor j/2\rfloor+\lfloor k/2\rfloor+\lfloor l/2\rfloor\not\equiv0\pmod{2}$,
so there are $4^4/4=64$ nonzero $t_{ijkl}$ variables.

Let $[[x_1,y_1,z_1],[x_2,y_2,z_2],[x_3,y_3,z_3],[x_4,y_4,z_4]]$ denote the form of $R_1$ and~$R_2$
as defined in Section~\ref{sec:background}.
Now count the number of points of type $00\mathord{*}\mathord{*}$ where the symbols `$*$' are arbitrary.
Note there are exactly $x_1x_2$ points which lie on both a line with index in $[0,x_1)$ and a line with index in $[10,10+x_2)$,
so
\begin{equation} \sum_{0\leq k,l\leq 3} t_{00kl} = x_1x_2 . \end{equation}
Similar equations can be derived by counting points of other types.
For example, there are exactly $x_1x_3$ points of type $0\mathord{*}0\mathord{*}$, exactly $y_1y_4$
points of type $1\mathord{*}\mathord{*}1$, and exactly $z_1x_2$ points of type $20\mathord{*}\mathord{*}$,
resulting in the equations
\begin{equation} \sum_{0\leq j,l\leq 3} t_{0j0l} = x_1x_3, \quad \sum_{0\leq j,k\leq 3} t_{1jk1} = y_1y_4, \quad \sum_{0\leq k,l\leq 3} t_{20kl} = z_1x_2 . \end{equation}

In case~3, the linear system corresponding to the $3^2\binom{4}{2}=54$ ways of fixing
two entries in the point type to values in $\{0,1,2\}$ when converted into reduced row echelon form
has a row corresponding to $t_{0123}-t_{3210}=-1/2$ which has
no integer solutions.

In cases~2 and~4, the linear systems corresponding to the $4^2\binom{4}{2}=96$ ways of fixing
two entries in the point type to values in $\{0,1,2,3\}$ are both inconsistent over the reals. \qed
\end{proof}

\end{document}
